% !TeX TS-program = lualatex
% !TeX root = main.tex
\DocumentMetadata{
    pdfversion=2.0,
    lang=it-IT,% default value: en-US
    % testphase={phase-III, table}% experimental, not working
}
\documentclass[
  % lingua del frontespizio
  language = english,% default: italian, accepted values: italian/english
  % debug,
  11pt
]{unitnthesis}

% todo: translates default strings if english is set

\unitnset{
    % valori accettati: c3a/cibio/cimec/cismed/dem/df/dfg/dicam/dii/dipsco/disi/dlf/dm/dsrs/sis
    department = disi,% default: disi, case insensitive field
    % valori accettati: bachelor/master/doctorate
    degree = bachelor,% default: bachelor, case insensitive field
    course = Informatica,
    title = Come scrivere la tesi in \LaTeX,
    subtitle = Il mio sottotitolo,
    keywords = { stem, algorithms, graphs },% use curly braces to delimit keywords
    % candidate/label = Candidata,
    candidate = Emanuele Nardi,
    candidate/ID = 000000,
    advisor/label = Relatrice,
    advisor = Prof.\ Alberto Montresor,% '\ ' to make a shorter space character
    coadvisors/label = Correlatori,
    coadvisors = { Prof.\ Paola Quaglia,Prof.\ Fausto Giunchiglia, Prof.\ Marco Patrigniani },% attenzione non mettere spazi
    year = 2024/2025,
    dedication = {A tutti gli studenti dell'UniTN\\per i loro primi passi in \LaTeX},
    chapther style = compact,
    page limit = 15,% should I consider the bibliograhy in the count?
}

% todos
% ask opinion on tex chat (?)

% \unitnset{
    % numero di loghi e posizione
    % logo ~ path / first  = assets/images/logo/eit-digital-logo.png,
    % logo ~ path / second = assets/images/logo/marchio-unitrento-colore-it.png,
    % logo ~ path / third  = assets/images/logo/stockholm-university-logo.png,
% }

% alternative
% \pathlogoset{
    % eit       = assets/images/logo/eit-digital-logo.png,
    % unitn     = assets/images/logo/marchio-unitrento-colore-it.png,
    % stockholm = assets/images/logo/stockholm-university-logo.png,
    % logo order = { eit, unitn, stockholm },
% }

% interface3 useful functions
% \file_if_exist:nTF
% \file_if_exist_input:nF

% seleziona la lingua principale del documento main=...
\usepackage[main=italian, english]{babel}
% vengono importati altri pacchetti che potrebbero essere utili
% consulta: </unitnthesis/pacchetti link>
\usepackage{addiotional-packages}

\begin{document}
\frontmatter

    \maketitle
        
    % \begin{acknowledgements}
    %     Here you can write your acknowledgments.
        
    %     \kant
    % \end{acknowledgements}

    % \tableofcontents
    % \listoffigures% (optional)
    % \listoftables% (optional)
    % \cleardoublepage

% \mainmatter

    % \begin{abstract}
    %     The \texttt{unitnthesis} class produces a template for the Bachelor, Master and Doctoral thesis manuscripts for students at UniTN.
    %     The class allows to write your manuscript using either italian or english; so, if you are an international student, feel free to use it!
    % \end{abstract}

    % \chapter{La classe \texttt{unitnthesis}}
% \chapter[Shorter version]{Title can be longer so a shorter version can be provided}
% [shorter version] is shown both in the header of pages (odd pages, top right) and in the index.

La classe \texttt{unitnthesis} è stata progettata per consentire allo studente di concentrarsi sul contenuto della tesi, piuttosto che sulla sua forma, in linea con la 
filosofia di \LaTeX.

\section{Opzioni del documento}

Utilizzando l'opzione \enquote{\texttt{code}} nel documento, verranno automaticamente caricati i pacchetti \package{minted} e \package{algorithm2e}. Il pacchetto \package{minted} consente di evidenziare la sintassi del codice in vari linguaggi di programmazione, mentre \package{algorithm2e} permette di scrivere algoritmi in modo formattato e leggibile.

\section{Pacchetti inclusi in questa classe}

Ecco la lista dei pacchetti con i link ai rispettivi siti di CTAN:

\begin{itemize}
    \item \package{babel} per il supporto multilingua;
    \item \package{xcolor} per la gestione dei colori;
    \item \package{csquotes} per la gestione delle virgolette (stile \texttt{italian=quotes});
    \item \package{microtype} per la tipografia;
    \item \package{geometry} per la configurazione dei margini;
    \item \package{setspace} per la gestione dell'interlinea nel testo;
    \item \package{parskip} per la gestione dell'interlinea;
    \item \package{titlesec} per lo stile dei titoli;
    \item \package{tocbibind} per aggiungere automaticamente la bibliografia e/o l'indice al sommario;
    \item \package{fancyhdr} per la gestione di intestazioni e piè di pagina;
    \item \package{enumitem} per l'utilizzo avanzato delle liste (opzione \texttt{inline});
    \item \package{graphicx} per la gestione delle immagini;
    \item \package{wrapfig2} per inserire immagini e tabelle a bordo del testo;
    \item \package{tabularx} per la gestione delle tabelle;
    \item \package{booktabs} per l'uso di filetti nelle tabelle;
    \item \package{tabularray} per tabelle avanzate;
    \item \package{tblr-extras} per tabelle avanzate che vanno su più pagine;
    \item \package{caption} per le didascalie delle immagini e tabelle;
    \item \package{subcaption} per le didascalie delle sotto-immagini;
    \item \package{footmisc} per la tipografia delle note a piè di pagina;
    \item \package{biblatex} per la gestione della bibliografia (con \texttt{backend=biber});
    \item \package{hyperref} per la gestione dei link cliccabili;
    \item \package{cleveref} per l'utilizzo delle etichette intelligenti.
\end{itemize}

Per coloro che frequentano corsi scientifici che richiedono la scrittura di codice è possibile caricare i pacchetti:
\begin{itemize}
    \item \package{algorithm2e} per la scrittura degli algoritmi;
    \item \package{minted} per l'inserimento di codice.
\end{itemize}

\section{Struttura delle cartelle}

La cartella \texttt{assets/} contiene diverse sottocartelle, ognuna con una funzione specifica per organizzare i diversi tipi di contenuto nella tesi:

\begin{enumerate}
    \item \texttt{chapters}: per i file di ciascun capitolo del documento;
    \item \texttt{appendixes}: per gli eventuali appendici aggiuntivi alla tesi;
    \item \texttt{images}: per le figure utilizzate nel documento, come fotografie \texttt{.jpg} o illustrazioni \texttt{.png};
    \item \texttt{resources}: per altre risorse utili, come dati, file di configurazione o materiali di ricerca o il database bigliografico \texttt{.bib} gestito dal pacchetto \package{biblatex};
    % \item \texttt{algorithms}: per memorizzare file relativi agli algoritmi e codici sorgente.
    % \item \texttt{codes}: per conservare il codice sorgente o esempi di codice, specialmente se sono inclusi tramite il pacchetto \package{minted};
    % \item \texttt{diagrams}: per archiviare diagrammi o grafici, come file \texttt{.tex} o \texttt{.pdf} prodotti dal pacchetto \package{tikz};
    % \item \texttt{tables}: per file relativi a tabelle e dati numerici, come file \texttt{.csv} o \texttt{.tex}.
\end{enumerate} 

Questa organizzazione aiuta a mantenere il progetto della tesi ben strutturato e facilmente gestibile.
Sentiti libero di eliminare qualsiasi cartella che non ti risulta utile.

Saranno mostrati diversi esempi per gestire questa struttura.
    % \chapter{Second Chapter}% level 0

Numbering of the sections is performed automatically by LaTeX, so don't bother adding them explicitly, just insert the heading you want between the curly braces.
If you don't want sections number, then add an asterisk \texttt{(*)} after the section command, but before the first curly brace, e.g.,

\verb|\section*{A Title Without Numbers}|.

\section{section}% level 1

\subsection{subsection}% level 2

\subsubsection{subsubsection}% level 3

\paragraph{paragraph}% level 4

\subparagraph{subparagraph}% level 5
    % \input{assets/chapters/03-development}
    % \input{assets/chapters/04-development}
    % \input{assets/chapters/05-conclusions}

    % \begin{appendices}
    %     \chapter{Qualche esempio}

\section{Come scrivere il testo.}

Vai a capo dopo ogni punto.
E qui scrivi la frase successiva.

Vai a capo due volte se vuoi iniziare un nuovo paragrafo.

Try to \emph{empathise words} or make them \textbf{bold}.

\section{Example lists}

Non numerated lists:
\begin{itemize}
  \item list entries start with the \texttt{item} command;
  \item individual entries are indicated with a black dot, a so-called bullet;
  \item the text in the entries may be of any length.
\end{itemize}

Numerated list:
\begin{enumerate}
  \item items are numbered automatically;
  \item the numbers start at 1 with each use of the \texttt{enumerate} environment;
  \item another entry in the list.
\end{enumerate}

Una lista nel testo:
\begin{itemize*}
    \item una lista in linea;
    \item un altro elemento;
    \item un altro ancora.
\end{itemize*}

Una lista con lettere maiuscole:
\begin{enumerate*}[label=\Alph*)]
    \item una lista in linea;
    \item un altro elemento;
    \item un altro ancora.
\end{enumerate*}

Una lista con lettere minuscole:
\begin{enumerate*}[label=\alph*.]
    \item una lista in linea;
    \item un altro elemento;
    \item un altro ancora.
\end{enumerate*}

\section{Example figure and table}

Una semplice tabella inserita nel testo.
Il testo, a meno che non si tratti di intestazioni, dovrebbe essere sempre allineato a sinistra.
Evita di utilizzare le linee verticali ed utilizza solo quelle orizzontali.

% @{} c c @{} indica come sono strutturate le colonne
% @{}: rimuove gli spazi sul bordo della tabella
% l (elle): indica la presenza di una colonna nel testo il cui testo è centrato
% puoi allineare il testo al centro con la colonna 'c', o a destra con la colonna 'r'
\begin{tabular}{@{} l l @{}}
\toprule
    intestazione & intestazione \\
\midrule
    corpo & corpo \\
\bottomrule
\end{tabular}

% Posizione nello specificatore float facoltativo
% h: here (qui)
% t: top (all'inizio della pagina)
% b: bottom (alla fine della pagina)
% p: page (una pagina dedicata a questa tabella)
\begin{table}[h]
    \centering
    % inserisci la didascalia *sopra* la tabella e non lasciare uno spazio fra 'caption' e 'label'
    \caption{Caption}\label{tab:my_label}
    \begin{tabular}{@{} *{2}{l} @{}}
    \toprule
        intestazione & intestazione \\
    \midrule
        corpo & corpo \\
    \bottomrule
    \end{tabular}
\end{table}

\section{Example figure and table with the tabularray package}

Attenzione il pacchetto \package{tabularray} seppur molto consuma prende molto tempo di compilazione, è sconsigliato utilizzarlo su Overleaf.
Per questo motivo è presente il codice \verb|\IfPackageLoadedTF{tabularray}| il quale controlla se il paccheto è stato caricato e in caso affermativo esegue il codice, altrimento stampa semplicemente un avvertimento.

\IfPackageLoadedTF{tabularray}{
La tabella~\ref{tab:tabella-con-tabularray}

\begin{table}[h!tbp]% here top bottom page
    \centering% allineata al centro
    \caption{Caption}% <- importante, non rimuovere questo '%'
    \label{tab:tabella-con-tabularray}% tab convenzione per 'table'
    \begin{tblr}{
    % opzioni della tabella
        % colspec={ *{3}{Q[c,m]} },% 3 colonne, il cui contenuto è centrato orizzontalmente e verticalmente
        row{1}={font=\bfseries},
        hline{1,Z} = {2pt},% prima ed ultima riga
        hline{2} = {1pt},
        hline{3-Y} = {.1pt, lightgray},% Y penultima riga
    }
    % contenuto della tabella
    Alpha & Beta & Gamma & Delta \\
    Epsilon & Zeta & Eta & Theta \\
    Iota & Kappa & Lambda & Mu \\
    Nu & Xi & Omicron & Pi \\
    Rho & Sigma & Tau & Upsilon \\
    Phi & Chi & Psi & Omega \\
    \end{tblr}
\end{table}

\begin{tblr}{
    colspec={ Q[l] Q[c] Q[r] },% allineamento colonne: (l)eft, (c)enter, (r)ight
    rowspec={ Q[t] Q[m] Q[b] },% allineamento right: (t)op, (m)iddle, (b)ottom
    hlines, vlines,% for debug purposes only
}
    {Alpha \\ Alpha} & Beta & Gamma \\% puoi spezzare il contenuto di una cella usando le '{}'
     Delta & Epsilon & {Zeta \\ Zeta} \\
     Eta & {Theta \\ Theta} & Iota \\
\end{tblr}

Esempio di tabella che prende tutta la pagina:

\begin{tblr}{
    colspec={@{} XXX @{}},
    hline{1,Z},% prima ed ultima riga
    vlines,% for debug purposes only
}
    Alpha & Beta & Delta\\
    Gamma & Delta & Epsilon\\
    Zeta & Eta & Theta\\
\end{tblr}

\begin{tblr}{
    hlines = {1pt,white},
    row{odd} = {blue7},
    row{even} = {azure7},
    column{1} = {purple7,c},
}
    Alpha & Beta & Gamma & Delta \\
    Epsilon & Zeta & Eta & Theta \\
    Iota & Kappa & Lambda & Mu \\
    Nu & Xi & Omicron & Pi \\
    Rho & Sigma & Tau & Upsilon \\
    Phi & Chi & Psi & Omega \\
\end{tblr}

\begin{tblr}{
    row{odd} = {bg=lightgray},
    % l'ordine dei comandi è importante
    row{1} = {bg=darkgray, fg=white, font=\sffamily},
}
    Alpha & Beta & Gamma \\
    Delta & Epsilon & Zeta \\
    Eta & Theta & Iota \\
    Kappa & Lambda & Mu \\
    Nu Xi Omicron & Pi Rho Sigma & Tau Upsilon Phi \\
\end{tblr}

\(\begin{tblr}{
    column{1} = {mode=text},
    column{3} = {mode=dmath},% display math
}
    Alpha   & \frac{1}{2} & \frac{1}{2} \\
    Epsilon & \frac{3}{4} & \frac{3}{4} \\
    Iota    & \frac{5}{6} & \frac{5}{6} \\
\end{tblr}\)
}{Il pacchetto \texttt{tabularray} non è stato caricato.}

\section{Wrap text of tables in text}

\IfPackageLoadedTF{wrapfig2}{
\begin{wrapfigure}{r}{50mm}
\centering\unitlength=1mm
\begin{picture}(40,30)
\polyline(0,0)(0,30)(40,0)(0,0)(40,30)(0,30)
\Line(40,0)(40,30)
\end{picture}
\caption{A rectangle with its diagonals}\label{fig:figure}
\end{wrapfigure}
{\itshape \kant[1]}
}{Il pacchetto \texttt{wrapfig2} non è stato caricato.}

For more example take a look at the package documentation of \package{wrapfig2}.

\section{Inserimento del codice}

Attenzione il pacchetto \package{minted} seppur molto consuma prende molto tempo di compilazione, è sconsigliato utilizzarlo su Overleaf.
Per questo motivo è presente il codice \verb|\IfPackageLoadedTF{minted}| il quale controlla se il paccheto è stato caricato e in caso affermativo esegue il codice, altrimento stampa semplicemente un avvertimento.

% \IfPackageLoadedTF{minted}{
% Inserimento di uno spezzone breve di codice direttamente nel testo:
% \begin{minted}{java}
% public class Hello {
%    public static void main(String[] args) {
%       System.out.println("Hello, World!");
%    }
% }
% \end{minted}
% }{Il pacchetto \texttt{minted} non è stato caricato.}
    %     \chapter{Department guidelines}

% \section{c3a}
% \section{cibio}
% \section{cimec}
% \section{cismed}
% \section{dem}
% \section{df}
% \section{dfg}
% \section{dicam}
% \section{dii}
% \section{dipsco}
% \section{disi}
% \section{dlf}
% \section{dm}

\section{Dipartimento di Sociologia e Ricerca Sociale}

All'interno del documento \enquote{\emph{\href{https://infostudenti.unitn.it/it/conseguimento-titolo-corsi-di-laurea-sociologia-e-ricerca-sociale}{linee guida per la redazione dell'elaborato finale dl laurea e della tesi dl laurea magistrale}}}, nella sezione \enquote{consigli redazionali} viene riportato:

\subsection{Formato dell'elaborato finale della laurea triennale}

L'elaborato scritto deve essere lungo all'incirca circa \num{90000} caratteri spazi inclusi (circa \num{13000} parole), carattere Arial o Helvetica; corpo del del testo 12, delle note 10.
La tesi deve essere composta a interlinea \num{1.5} con ampi margini laterali (3cm almeno per ogni lato).

\subsection{Struttura dell'elaborato finale e della tesi magistrale}

Generalmente una tesi di laurea magistrale si compone di una breve introduzione, di quattro/cinque capitoli e delle conclusioni.
La struttura dell'elaborato finale è sostanzialmente identica, anche se non composta da \enquote{capitoli} ma da \enquote{paragrafi}.
\begin{itemize}
    \item nell'\textbf{introduzione} vengono presentati i motivi di interesse della tesi e si delinea il modo in cui verrà svolto il tema;
    \item nel \textbf{capitolo 1} (o, per la laurea triennale, nel paragrafo 1) si passa in rassegna la letteratura rilevante sul tema in questione e vi si inquadra il tema scelto (la letteratura disponibile, gli antefatti storici, i precedenti studi sull'argomento, e cosi via), evidenziando, se del caso, l'originalità del lavoro svolto;
    \item nel \textbf{capitolo 2} (o, per la laurea triennale, nel paragrafo 2) si introduce la metodologia utilizzata;
    \item nei \textbf{due o tre restanti capitoli} (o, per la laurea triennale, nei paragrafi) si presentano i risultati ottenuti dalla propria ricerca (generalmente passando dal generale al particolare, dalle applicazioni più semplici a quelle più sofisticate, dal passato al presente e al futuro).
    \item nella \textbf{conclusione} si presentano i propri giudizi sui risultati ottenuti, sulla rispondenza della metodologia adottata allo scopo della ricerca, sui possibili sviluppi della ricerca.
\end{itemize}

La stesura dell'elaborato finale o della tesi viene svolta sulla base della scaletta condivisa con il/la supervisore o relatore/relatrice cominciando di solito dal Capitolo 1. Non c'è un ordine prefissato, ma è consigliabile che l'introduzione e le considerazioni conclusive vengano redatte per ultime: i loro contenuti sono particolarmente importanti: l'interesse del/della lettore/lettrice tende infatti a concentrarsi in particolare su di esse. Dal momento che riassumono l'intero lavoro, è opportuno
che il resto sia già scritto.

Chiaramente, sia l'elaborato finale che la tesi dovranno anche presentare un frontespizio, un indice iniziale e una bibliografia/sitografia finale.

% \section{Scuola di Studi Internazionali}
    % \end{appendices}

% \backmatter

    % Print the entire bibliography regardless of whether they have been cited in the text
    % \nocite{*}

    % Add contents to ToC
    % \addcontentsline{toc}{chapter}{Bibliography}% not needed anymore, to check
    
    % \printbibliography
\end{document}